\scanpage{271}
\mspage{12}

\uline{Corollaire} L'ensemble des \deleted{\unsure{points}}\unsure{sommets} critiques de $I$ \deleted{\unsure{déf}} \unsure{est|peut} aussi l'ens. des sommets du \deleted{\unsure{la}} plus petit sous-\deleted{catégorie pleine \uline{connexe}} \inserted{intervalle} de $I$ contenant les extrémités.

Plus généralement, si $i<j$ sont deux sommets de $I$, le plus petit \inserted{ss-} \deleted{catégorie pleine} intervalle de $I$ contenant $i,j$ est formé de $i,j$ et des sommets critiques $s$ tels que $i<s<j$.

\uline{Définition} Un intervalle est dit \uline{réduit} \inserted{(ou \uline{\unsure{irrédondant|irrédundant}})} si tous ses sommets sont critiques, ou ce qui revient au même, si ses chaînes structurales sont de longueur 1.

Un intervalle orienté \uline{réduit} est connu, à isom unique près, quand on connaît

a) Sa longueur $n$

b), Sa direction initiale $+$ (directe) ou $-$ (rétrograde)

\uline{NB} Cette direction est la $\mathord{=}$ que que la direction finale, \deleted{\unsure{Dans le cas d}} si $n$ est impaire, elle est opposée si $n$ est paire.

Ainsi, suivant les deux cas \inserted{de direction initiale}, on trouve

\diagram{\begin{tikzcd}[ampersand replacement=\&,column sep=small]
0\arrow[r] \& 1 \& 2\arrow[l]\arrow[r] \& \ldots \& \ldots\,\unsure{\bullet|\cdot}\, n \& \left(\begin{gathered}\text{Avec }n-1\longrightarrow n\quad\text{si }n\text{ imp.}\\\phantom{\text{Avec }}n-1\longleftarrow n\quad\text{si }n\text{ pair}\end{gathered}\right)
\end{tikzcd}}

\diagram{\begin{tikzcd}[ampersand replacement=\&,column sep=small]
0 \& 1\arrow[l]\arrow[r] \& 2 \& \ldots\arrow[l] \& \unsure{\to|\cdot} \& \ldots\,n \& \left(\begin{gathered}\text{Avec }n-1\longleftarrow n\quad\text{si }n\text{ imp}\\\phantom{\text{Avec }}n-1\longrightarrow n\quad\text{si }n\text{ pair}\end{gathered}\right)
\end{tikzcd}}

\deleted{\uline{\unsure{Définition}}} \inserted{Prop. 6} \deleted{Le \deleted{plus} sous-intervalle de $I$ formé des sommets critiques est \deleted{appelé l'intervalle réduit} \inserted{non seulement le plus petit} sous-intervalle contenant les extrémités, mais aussi le plus grand sous-intervalle irrédundant.}

\deleted{\unsure{Car :}} un sous-intervalle $J$, ayant $i$ et $j$ ($i<j$) ``comme'' plus petit et plus grand sommet, est \deleted{irrédundant} \inserted{réduit} sss \deleted{\unsure{les som}} tout \deleted{\unsure{$j\in$}} $s\in J$ distinct de $i$ et de $j$ est critique \unsure{\}|?}.

\uline{Cor.} La ss-cat. pleine de $I$ formée des sommets critiques est un sous-intervalle. On l'appelle l'intervalle réduit \uline{associé} à $I$, noté $I_{\mathrm{red}}$.

% archive sheet number: 270
